\chapter{Fitting measured data}\label{ch:fitting}

Measured profiles, point clouds from a coordinate measuring machine, simulation output or
sampled analytic functions are turned into NURBS by \emph{fitting}. \NF{} offers two families of
methods, both following Piegl \& Tiller, chapter~9:

\begin{itemize}
\item \textbf{Interpolation} passes exactly through every data point. It needs as many control
  points as data points and reproduces noise faithfully.
\item \textbf{Least-squares approximation} uses fewer control points than data points and
  minimises the sum of squared deviations at the data parameters; the first and the last point
  are always interpolated. It smooths the data and gives compact curves.
\end{itemize}

All fits produce \emph{non-rational} B-splines on the domain $[0,1]$ (curves) or $[0,1]^2$
(surfaces). Input points are NumPy arrays: shape $(n, 3)$ for curves, $(n_u, n_v, 3)$ for
surface grids.

\section{Data parameters}\label{sec:data-params}

Every fit first assigns a parameter $\bar u_k \in [0,1]$ to each data point. The choice matters:
it determines where along the curve each point ends up and therefore the shape between the
points. \apiref{data_parameters}\code{(points, method)} returns these parameters; the enumeration
\apiref{Parametrization} offers three methods (P\&T Eq.~9.3--9.6):

\begin{tabularx}{\linewidth}{@{}l>{\raggedright\arraybackslash}X@{}}
\toprule
Method & Parameter increments \\
\midrule
\code{Uniform} & equal: $\bar u_k = k/(n-1)$. Only suitable for evenly spaced data. \\
\code{ChordLength} & proportional to the distance between consecutive points,
  $\Delta\bar u_k \propto |\mathbf Q_k - \mathbf Q_{k-1}|$. Good for evenly curved data. \\
\code{Centripetal} & proportional to the square root of the distance,
  $\Delta\bar u_k \propto |\mathbf Q_k - \mathbf Q_{k-1}|^{1/2}$. Robust at sharp turns and
  uneven spacing; the \textbf{default} of every fit. \\
\bottomrule
\end{tabularx}

Consecutive identical points are rejected (\code{GeometryError}) for the non-uniform methods,
since they would give a zero parameter increment. The data parameters are also what you need to
measure the deviation of a fit \emph{at the data}: the distance between $\mathbf C(\bar u_k)$ and
$\mathbf Q_k$ is an upper bound of the point-to-curve distance.

\section{Curve interpolation}\label{sec:interpolate-curve}

\begin{lstlisting}[style=text]
interpolate_curve(points, degree=3, method=Parametrization.Centripetal,
                  end_derivatives=None) -> NurbsCurve
\end{lstlisting}

\apiref{interpolate_curve} computes the global interpolant (A9.1): a B-spline of the given
degree through all $n$ points with $n$ control points and knots obtained by averaging the
parameters (Eq.~9.8). It needs at least $\text{degree}+1$ points. \code{end_derivatives} is an
optional pair $(\mathbf C'(0), \mathbf C'(1))$ of first derivatives \emph{with respect to the
curve parameter on $[0,1]$} -- not unit tangents. A natural magnitude is the total chord length
of the data. Prescribed end derivatives add two control points (Eq.~9.21/9.22) and require a
degree of at least 2.

Interpolation reproduces polynomial data of degree $\le p$ exactly, e.g.\ points on a line give a
line for any degree (with a suitable parametrisation). Between the points the interpolant may
oscillate if the data are noisy or unevenly spaced; check it on a dense sample.

\section{Least-squares approximation}\label{sec:approximate-curve}

\begin{lstlisting}[style=text]
approximate_curve(points, degree, n_ctrl, method=Parametrization.Centripetal) -> NurbsCurve
approximate_curve_tol(points, degree, tolerance, method=Parametrization.Centripetal) -> NurbsCurve
\end{lstlisting}

\apiref{approximate_curve} fits a B-spline with exactly \code{n_ctrl} control points
($\text{degree}+1 \le n_\text{ctrl} \le n$) in the least-squares sense (Eq.~9.63--9.69). The
interior knots are placed so that every knot span contains data; the first and last data point
are interpolated.

\apiref{approximate_curve_tol} searches the \emph{smallest} number of control points for which
every data point lies within \code{tolerance} of the curve point at its parameter: exponential
search followed by bisection, falling back to interpolation if no approximation meets the
tolerance. This is usually the most convenient fitting function -- specify the accuracy, not the
number of control points.

\section{Example: fitting a spiral}

\example{07_curve_fitting.py}

Output:

\begin{lstlisting}[style=text]
Uniform      parameters of points 0..3: [0.0, 0.0256, 0.0513, 0.0769]
ChordLength  parameters of points 0..3: [0.0, 0.0041, 0.0116, 0.0213]
Centripetal  parameters of points 0..3: [0.0, 0.0105, 0.0246, 0.0407]
interpolate_curve:       40 ctrl, data deviation 3.7e-16, distance to spiral 1.5e-05
approximate_curve n= 6:   6 ctrl, data deviation 1.9e-02
approximate_curve n= 8:   8 ctrl, data deviation 2.8e-03
approximate_curve n=12:  12 ctrl, data deviation 7.1e-04
approximate_curve_tol 1e-03:  12 ctrl, data deviation 7.1e-04
approximate_curve_tol 1e-05:  25 ctrl, data deviation 6.0e-06
with end derivatives: 42 ctrl (2 more), C'(0) = Vec3(0.000000, 4.000000, 0.000000)
line data: C(0.5) = Vec3(2.500000, 5.000000, 0.000000)
\end{lstlisting}

Observations:

\begin{itemize}
\item The data are denser at the start (the samples are spaced like $t^{1.5}$). Uniform
  parameters ignore this; chord length and centripetal parameters follow the actual spacing.
\item The interpolant is exact at the 40 data points ($3.7\cdot10^{-16}$), but between them it
  deviates from the true spiral by $1.5\cdot10^{-5}$ where the samples are sparse. The data
  deviation of a fit is \emph{not} its deviation from the underlying shape.
\item With 12 control points the least-squares curve stays within $7.1\cdot10^{-4}$ of the data.
  \code{approximate_curve_tol} with a tolerance of $10^{-3}$ finds exactly these 12 control
  points; a tolerance of $10^{-5}$ needs 25.
\end{itemize}

\section{Best-fit lines and planes}\label{sec:fit-line-plane}

Measured points often describe a straight edge or a flat face. \apiref{fit_line} and
\apiref{fit_plane} fit them by \emph{orthogonal} regression: they minimise the sum of the
squared perpendicular distances, unlike a parametric least-squares fit $z = f(x, y)$, which
measures errors along one axis and depends on the orientation of the coordinate system.

\begin{lstlisting}
lf = nf.fit_line(points)    # LineFit: point (centroid), direction, max_deviation, rms
pf = nf.fit_plane(points)   # PlaneFit: plane (through the centroid), max_deviation, rms
\end{lstlisting}

The direction of the line is the eigenvector of the largest eigenvalue of the scatter matrix of
the points, the plane normal that of the smallest one. Both are unit vectors whose
largest-magnitude component is positive, so the result does not depend on the order of the
points. \code{max_deviation} and \code{rms} report how well the points fit, e.g. as a flatness
check of a measured face. A line needs two distinct points, a plane three points that are not
collinear. \Cref{sec:project} has an example.

\section{Surface fitting}\label{sec:surface-fitting}

\begin{lstlisting}[style=text]
interpolate_surface(points, degree_u=3, degree_v=3, method=Parametrization.Centripetal) -> NurbsSurface
approximate_surface(points, degree_u, degree_v, ctrl_u, ctrl_v,
                    method=Parametrization.Centripetal) -> NurbsSurface
\end{lstlisting}

Surface fits work on a \emph{grid} of points, a NumPy array of shape $(n_u, n_v, 3)$: \code{points[i, j]}
is the point at the $i$-th position along $u$ and the $j$-th along $v$ (build it with
\code{np.meshgrid(..., indexing="ij")}). For unordered points use \apiref{fit_surface}
(\cref{sec:scattered-fitting}).

\apiref{interpolate_surface} interpolates all grid points (A9.4) with an $n_u\times n_v$ control
net. The $u$ parameters are averaged over all columns and the $v$ parameters over all rows;
degenerate rows (all points equal, e.g.\ a pole) are skipped in the averaging.

\apiref{approximate_surface} fits a \code{ctrl_u} $\times$ \code{ctrl_v} net by least squares,
first through all rows and then through the resulting columns. The corner points are
interpolated.

For a graph $z = f(x,y)$ over a rectangle whose $u$ and $v$ should coincide with $x$ and $y$, use
\apiref{sag_surface} (\cref{ch:optics}) instead: it keeps $x$ and $y$ exactly linear and
refines to a tolerance.

\example{07_surface_fitting.py}

Output:

\begin{lstlisting}[style=text]
grid shape (12, 10, 3)
interpolate_surface: net 12 x 10, max distance to z = sin(x) cos(y) 3.0e-05, corner error 5.0e-16
approximate_surface: net 6 x 5, max distance to z = sin(x) cos(y) 6.5e-04, corner error 0.0e+00
interpolant: max distance of the 120 grid points 3.8e-14
\end{lstlisting}

\section{Fitting scattered points}\label{sec:scattered-fitting}

\begin{lstlisting}[style=text]
fit_surface(points, num_u=0, num_v=0, degree_u=3, degree_v=3, smoothing=1e-3,
            iterations=2, base=None, boundary=[]) -> FittedSurface
\end{lstlisting}

\apiref{fit_surface} fits a non-rational B-spline surface (uniform clamped knots, domain
$[0,1]^2$) to \emph{unordered} points, an array of shape $(n, 3)$, e.g.\ the measurement of an
optical freeform surface or a topography. It minimises
\[
  \frac1n\sum_k \bigl|\mathbf S(u_k, v_k) - \mathbf P_k\bigr|^2
  + \lambda \iint |\mathbf S_{uu}|^2 + 2|\mathbf S_{uv}|^2 + |\mathbf S_{vv}|^2 \,du\,dv ,
\]
the mean squared deviation plus a thin-plate energy, as one banded linear system for the
control points. The weight $\lambda$ is \code{smoothing} times the ratio of the traces of the
two matrices, so the value is relative and does not depend on the size of the data or on the
number of points. It also fills regions without data smoothly (the corners of the parameter
rectangle over a round aperture in the example). Plane-like data are reproduced exactly for
any smoothing, since planes have no bending energy.

The point parameters $(u_k, v_k)$ come from
\begin{itemize}
  \item the best-fit plane of the points (default): their coordinates in the plane, aligned with
    the principal direction of the points and scaled to the bounding rectangle. This suits
    height fields over a plane;
  \item \code{base}, a surface close to the data: the closest points on it. Use this for data
    that fold over the best-fit plane, such as a cylinder segment;
  \item \code{boundary}, four curves \code{[c0, c1, d0, d1]} as for \code{coons}: the Coons
    patch of the curves is the base, and the fitted surface \emph{interpolates the curves
    exactly} (its boundary control rows are those of the refined patch; the curves must be
    non-rational).
\end{itemize}
After each fit the points are projected onto the surface (\code{closest_point_local}) and the
surface is fitted again with the corrected parameters, \code{iterations} times. The result
\apiref{FittedSurface} holds the \code{surface}, the largest and rms distances of the points from
it (\code{max_error}, \code{rms_error}) and the foot point \code{parameters} as an $(n, 2)$
array.

The control net has \code{num_u} $\times$ \code{num_v} points (0: about $\sqrt{n/4}$ per
direction, at most 40). For smooth, exact data keep the default \code{smoothing}: $10^{-3}$ barely
changes the fit. For noisy data, use fewer control points and/or more smoothing (0.1 to 1),
and compare the rms residual with the expected noise level: a residual far below it means the
surface follows the noise.

\example{07_scattered_fitting.py}

Output:

\begin{lstlisting}[style=text]
exact data: 27 x 27 control points, max error 2.2e-05, rms 2.1e-06
  deviation from f between the points (inner disc): 2.8e-06
noisy data, 27 x 27, smoothing 0.001: rms residual 1.76e-03, deviation from f 3.2e-03
noisy data, 27 x 27, smoothing 1: rms residual 1.83e-03, deviation from f 1.7e-03
noisy data, 12 x 12, smoothing 0.1: rms residual 1.93e-03, deviation from f 7.2e-04
boundary curves: max error 1.2e-05, rim height 0.0e+00
half cylinder with base: max error 8.4e-05
\end{lstlisting}

With noise $\sigma = 0.002$, the residuals stay near $\sigma$ in all three fits, but only the
smoothed and coarser fits are closer to the true surface than the noise itself.

\section{Practical advice}\label{sec:fitting-advice}

\begin{itemize}
\item \textbf{Degree.} Cubic (3) is the standard: $C^2$ continuous, flexible, little overshoot.
  Use degree 5 for very smooth data where curvature continuity of higher derivatives matters;
  degree 1 or 2 only for polylines or special cases.
\item \textbf{Interpolate or approximate?} Interpolate exact data (analytic samples, design
  points). Approximate measured data: noise is smoothed, and the curve has far fewer control
  points, which makes all later operations faster.
\item \textbf{Prefer \code{approximate_curve_tol}.} A tolerance a little above the measurement
  uncertainty gives the most compact curve that still represents the data.
\item \textbf{Check between the points.} The tolerance is verified at the data parameters. If the
  shape between the samples matters, evaluate the fit on a denser set of points (as in the
  examples, with \apiref{closest_point}) or sample the source more densely.
\item \textbf{Sharp corners.} A single smooth fit through a corner rounds it off and overshoots.
  Split the data at the corner, fit the pieces separately and \apiref{join} them.
\item \textbf{Closed curves.} The fits produce open, clamped curves. For a closed curve repeat
  the first point at the end; the curve is then closed but in general only $C^0$ at the seam.
  Prescribe equal end derivatives with \code{end_derivatives} to make it $C^1$.
\end{itemize}
